%! From Combinations to Solutions --- Setting Up Ax = b — Unit 2, Lesson 2.0 — Group Activity (Answer Key)
\documentclass[10pt]{article}
\usepackage{linalg-article}
\usepackage{linalg-key}   % pulls in -boxes; do NOT also load -boxes
\newcommand{\TallMath}[1]{$\displaystyle #1\rule[-1.4em]{0pt}{3.2em}$}
\newcommand{\xx}{\mathbf{x}}
\newcommand{\bb}{\mathbf{b}}

\begin{document}

\pageheader{Unit 2, Lesson 2.0}{Group Activity --- Answer Key}

\vspace{0.10in}

\begin{headlinebox}{blush}
\small\textbf{Blend the Batch.} A smoothie bar mixes two base powders. Each scoop of a base adds
a fixed amount of \emph{protein} and \emph{fiber}. To hit a nutrition target you solve
$A\xx=\bb$: the bases are the \textbf{columns} of $A$, the scoops are $\xx$, the target is $\bb$.
Work with your partner and \emph{show your steps}.
\end{headlinebox}

\vspace{0.10in}

\begin{tcolorbox}[colback=white, colframe=black!40, title=\textbf{Tier R --- Remediate},
    fonttitle=\bfseries, arc=1.5mm, left=3mm, right=3mm, top=2mm, bottom=2mm, breakable]
Solve $A\xx=\bb$ with $A=\begin{bmatrix} 1 & 1 \\ 1 & -1 \end{bmatrix}$ and
$\bb=\begin{bmatrix} 7 \\ 1 \end{bmatrix}$.
\begin{enumerate}[leftmargin=1.7em, itemsep=8pt]
  \item Write the two equations (one per row). \ansline{$x_1+x_2=7$;\ \ $x_1-x_2=1$}
  \item \emph{Add} the two equations to eliminate $x_2$, and find $x_1$.
\begin{work}
  (x_1+x_2)+(x_1-x_2) &= 7+1 \\
                 2x_1 &= 8 \\
                  x_1 &= 4
\end{work}
  \item Back-substitute to find $x_2$.
\begin{work}
  4 + x_2 &= 7 \\
      x_2 &= 3
\end{work}
  \item \textbf{Check:} rebuild $\bb$ as $x_1\begin{bsmallmatrix} 1\\1 \end{bsmallmatrix}+x_2\begin{bsmallmatrix} 1\\-1 \end{bsmallmatrix}$.
\begin{work}
  4\begin{bsmallmatrix} 1\\1 \end{bsmallmatrix} + 3\begin{bsmallmatrix} 1\\-1 \end{bsmallmatrix}
      &= \begin{bsmallmatrix} 7\\1 \end{bsmallmatrix} \ \ \text{\checkmark}
\end{work}
\end{enumerate}
\end{tcolorbox}

\begin{tcolorbox}[colback=white, colframe=black!40, title=\textbf{Tier A --- Approaching Proficiency},
    fonttitle=\bfseries, arc=1.5mm, left=3mm, right=3mm, top=2mm, bottom=2mm, breakable]
Each scoop of \textbf{Base X} adds $2$ g protein and $1$ g fiber; each scoop of \textbf{Base Y}
adds $1$ g protein and $3$ g fiber. A smoothie must total $8$ g protein and $9$ g fiber.
\begin{enumerate}[leftmargin=1.7em, itemsep=8pt]
  \item Set it up as $A\xx=\bb$. Fill in the matrix and target.
        \[ A=\begin{bmatrix} {\color{keyred}\mathbf{2}} & {\color{keyred}\mathbf{1}} \\[3pt]
                             {\color{keyred}\mathbf{1}} & {\color{keyred}\mathbf{3}} \end{bmatrix},\qquad
           \bb=\begin{bmatrix} {\color{keyred}\mathbf{8}} \\[3pt] {\color{keyred}\mathbf{9}} \end{bmatrix} \]
  \item Write the system and solve it by elimination.
\begin{work}
  2(x_1+3x_2)-(2x_1+x_2) &= 18-8 \\
                    5x_2 &= 10 \\
                     x_2 &= 2 \\
                     x_1 &= 3
\end{work}
  \item \emph{Interpret:} how many scoops of each base, and does the answer make sense in cups?
        \ansline{$3$ scoops of Base X and $2$ of Base Y --- both whole, positive scoops, so it is a valid recipe.}
  \tcbbreak
  \item A new target is $8$ g protein and $2$ g fiber. Solve again --- what goes wrong, and what
        would it mean to ``use $-\tfrac{4}{5}$ of a scoop''?
\begin{work}
  2(x_1+3x_2)-(2x_1+x_2) &= 4-8 \\
                    5x_2 &= -4 \\
                     x_2 &= -\tfrac{4}{5} \\
                     x_1 &= \tfrac{22}{5}
\end{work}
        \ansline{The algebra still has a solution, but $x_2$ is a \emph{negative} scoop ---
        impossible in real life, since you cannot remove Base Y you never added.}
\end{enumerate}
\end{tcolorbox}

\begin{tcolorbox}[colback=white, colframe=black!40, title=\textbf{Tier E --- Extension},
    fonttitle=\bfseries, arc=1.5mm, left=3mm, right=3mm, top=2mm, bottom=2mm, breakable]
Two systems have \emph{parallel} column directions. For each, decide how many solutions it has
and \textbf{justify} from both the picture (two lines) and reachability of $\bb$.
\begin{enumerate}[leftmargin=1.7em, itemsep=8pt]
  \item $\begin{aligned}[t] x_1 + 2x_2 &= 4 \\ 2x_1 + 4x_2 &= 9 \end{aligned}$
        \hspace{1.5em} How many solutions? Why?
\begin{work}
  2(x_1+2x_2) &= 2(4) = 8 \\
            8 &\ne 9
\end{work}
        \ansline{\textbf{None.} The lines are parallel but different. Both columns point along
        $\begin{bsmallmatrix} 1\\2 \end{bsmallmatrix}$, so $C(A)$ is that line and
        $\begin{bsmallmatrix} 4\\9 \end{bsmallmatrix}$ is off it --- $\bb$ is unreachable.}
  \item $\begin{aligned}[t] x_1 + 2x_2 &= 3 \\ 2x_1 + 4x_2 &= 6 \end{aligned}$
        \hspace{1.5em} How many solutions? Why?
\begin{work}
  2(x_1+2x_2) &= 2(3) = 6 \\
  \begin{bsmallmatrix} 3\\6 \end{bsmallmatrix} &= 3\begin{bsmallmatrix} 1\\2 \end{bsmallmatrix}
\end{work}
        \ansline{\textbf{Infinitely many.} Equation~2 is just twice equation~1 --- the same line.
        Here $\bb$ \emph{is} on the column line, so every point of that line solves it.}
\end{enumerate}
\end{tcolorbox}

\end{document}
